\clearpage
\part{Quantum Mechanics}
  \label{part:quantum-mechanics}

  This part presents a unified treatment of quantum phenomena as consequences of non-injective
  projection rather than as fundamental departures from classical logic.
  Quantum correlations, indeterminacy, and entanglement are derived from the structural
  non-factorizability of projected descriptions of the underlying substrate~$\chi$.

  Violations of Bell inequalities are interpreted as signatures of ontological projection
  failure, not as evidence of superluminal influences.
  Measurement, decoherence, and apparent wave-function collapse are described as
  context-dependent stabilization processes within projected regimes.

  Temporal ordering, established in Part~\ref{part:foundations} as the monotone growth
  of cumulative projected capacity $I(n)$ along the admissible cascade, provides the
  pre-existing ordering with respect to which effective quantum evolution is formulated.
  The Schr\"odinger equation is accordingly not a fundamental law from which time
  originates, but an effective description formulated within a structurally prior ordering.

  The relation to the standard quantum formalism is clarified, including the emergence of
  Hilbert space structure, effective Schr\"odinger dynamics, and the Born rule.
  Quantum mechanics is thus situated as a consistent and predictive effective framework
  arising from deeper relational constraints.

  The logical chain underlying this part can be summarised as follows.
  Non-injectivity of~$\Pi$ (established as a structural necessity in
  Section~\ref{subsec:non-injectivity-not-ad-hoc} and~\cite[Thm.~3]{Beau2026l}) implies that distinct configurations
  of~$\chi$ are mapped to the same observable: a given observable~$o$ carries
  an irreducible fiber $\Pi^{-1}(o)$ of indistinguishable pre-images.
  Because the effective description has no direct access to~$\chi$, it must represent
  the coexistence of projectively admissible contributions internally, without appeal
  to the substrate.
  This necessity forces an additive, phase-sensitive, norm-preserving structure on the
  space of effective descriptions---a complex Hilbert space, interpreted not as a
  fundamental ontology but as the minimal representational closure compatible with
  non-injective projection (Section~\ref{sec:relation-to-quantum-formalism}).
  In the regime where the quasi-stable background and Born--Infeld bounds jointly hold,
  this structure supports a unique effective dynamics of Schr\"odinger type, derived in
  Appendix~\ref{app:schrodinger-derivation} as Theorem~\ref{thm:schrodinger-emergence}.
  The spinorial structure $\mathrm{SU}(2)$, required by the maximum quaternionic
  admissibility of the projection fiber
  (Section~\ref{subsec:binary-polyhedral-maximality}), completes the identification
  of the standard quantum-mechanical framework as an emergent effective description.
  The first three steps of this chain are structural necessities; the Hilbert
  interpretation is a structurally motivated consequence, not yet a fully derived
  theorem; Schr\"odinger dynamics is a proved result under explicit hypotheses.

  \clearpage
  \input{4-quantum_mechanics/09-quantum_phenomena_and_entanglement/quantum_phenomena_and_entanglement}
